% SPDX-License-Identifier: Apache-2.0
\section{Two Events at One Edge}
\label{sec:x-join}

An adder needs a word from each of two channels, so its event is
both arrivals: \code{parallel!(a.wait(), b.wait()).await}, which
completes when both waits have, with the two words. Producer A offers
every cycle and producer B every other one. A's word is taken when it
arrives and held inside the group until B's arrives, A's channel holds
the next offer meanwhile, and so A runs at B's rate with no code for
the backpressure. Figure~\ref{fig:x-join} shows it: \code{a\_valid}
stays high through the cycle A waits, and a sum lands on every edge B
delivers.

\begin{figure}[htbp]
\centering
\input{join_timing.tex}
\caption{The join. A's offer is taken at once and held; B's arrives every other cycle; the sum follows B.}
\label{fig:x-join}
\end{figure}

\lstinputlisting[caption={\code{ex\_join.rs}. Run with \code{bazel run //lib/examples:ex\_join}.},label={lst:x-join}]{lib/examples/ex_join.rs}

\lstinputlisting[style=ascii,caption={What \code{ex\_join} printed when the build ran it.},label={lst:x-join-out}]{out_join.txt}
